%! AP Statistics — Unit 4, Lesson 4.5 — Guided Notes
\documentclass[10pt]{article}
\usepackage{apstats-article}
\usepackage{apstats-boxes}

\begin{document}

\pageheader{Unit 4, Lesson 4.5}{Guided Notes}
\namedateperiod

\vspace{0.12in}

\begin{objectivebox}
\small By the end of this lesson, I will be able to\dots\\[4pt]
\writeline\\[1pt]
\writeline\\[1pt]
\writeline
\end{objectivebox}

\vspace{0.1in}

\begin{vocabbox}
\termblanklong{Conditional probability $P(A|B)$}
\termblanklong{Conditioning event}
\termblanklong{Multiplication rule}
\termblanklong{Marginal probability}
\end{vocabbox}

\vspace{0.1in}

\begin{hookbox}
\small
\textbf{Card scenario:} I draw a card and tell you it is a face card.
\begin{enumerate}[leftmargin=1.5em, itemsep=4pt]
  \item Before I told you, what was $P(\text{King})$? \blank{4cm}
  \item Now that you know it's a face card, what is the new probability it's a King?
        \blank{4cm}
  \item How many face cards are there? \blank{2cm} \quad How many are Kings? \blank{2cm}
\end{enumerate}
\textbf{Why did the probability change?}\\[4pt]
\writeline
\end{hookbox}

\vspace{0.1in}

\begin{notesbox}{Conditional Probability (VAR-4.D.1)}
\small
\textbf{Definition:} The \textbf{conditional probability} $P(A|B)$ is the probability
that event $A$ will occur \blank{4cm} event $B$ \blank{3cm}.

\vspace{8pt}
\textbf{Formula:}
\[
  P(A|B) = \frac{P(A \cap B)}{P(B)}, \quad \text{provided } P(B) \blank{2cm}.
\]

\vspace{6pt}
\textbf{Interpretation:} Knowing $B$ occurred \blank{5cm} the sample space to only
outcomes in $B$.

\vspace{10pt}
\textbf{From a two-way table:}
\begin{enumerate}[leftmargin=1.5em, itemsep=4pt]
  \item Find $P(A \cap B)$ = \blank{5cm} $\div$ grand total
  \item Find $P(B)$ = \blank{5cm} $\div$ grand total
  \item Compute $P(A|B)$ = \blank{5cm} $\div$ \blank{3cm} (the grand total cancels!)
\end{enumerate}

\vspace{8pt}
\textbf{Short-cut for tables:} $P(A|B) = \dfrac{\text{count in both }A\text{ and }B}{\blank{4cm}}$
\end{notesbox}

\vspace{0.1in}

\begin{notesbox}{Worked Example --- Two-Way Table}
\small
Use the sport and training table below (120 athletes):
\vspace{4pt}
\renewcommand{\arraystretch}{1.4}
\begin{tabularx}{\linewidth}{@{} l X X r @{}}
 & \textbf{Trains Daily} & \textbf{Trains Weekly} & \textbf{Total} \\
\hline
\textbf{Runs}     & 24 & 16 & 40 \\
\textbf{Swims}    & 18 & 22 & 40 \\
\textbf{Cycles}   & 20 & 20 & 40 \\
\hline
\textbf{Total}    & 62 & 58 & 120 \\
\end{tabularx}

\vspace{8pt}
\textbf{Find $P(\text{Runs} | \text{Trains Daily})$:}
\begin{itemize}[itemsep=3pt]
  \item $P(\text{Runs} \cap \text{Trains Daily}) = $ \blank{3cm}
  \item $P(\text{Trains Daily}) = $ \blank{3cm}
  \item $P(\text{Runs} | \text{Trains Daily}) = \dfrac{\blank{2cm}}{\blank{2cm}} = $ \blank{3cm}
\end{itemize}

\vspace{6pt}
\textbf{Interpret in context:} Among athletes who train daily,
\blank{4cm}\% run.
\end{notesbox}

\vspace{0.1in}

\begin{notesbox}{Multiplication Rule (VAR-4.D.2)}
\small
\textbf{Rearranging the formula} gives the \textbf{multiplication rule}:
\[
  P(A \cap B) = P(A) \cdot P(B|A) = P(B) \cdot P(A|B)
\]

\vspace{6pt}
\textbf{Use this when:} you know $P(A)$ and $P(B|A)$ and need to find \blank{4cm}.

\vspace{8pt}
\textbf{Example:} $P(\text{Trains Daily}) = 62/120$. $P(\text{Runs}|\text{Trains Daily}) = 24/62$.
\[
  P(\text{Runs} \cap \text{Trains Daily}) = \blank{3cm} \times \blank{3cm} = \blank{3cm}
\]
\end{notesbox}

\vspace{0.1in}

\begin{practicebox}
\small Use the same 120-athlete table above.

\vspace{6pt}
\begin{enumerate}[leftmargin=1.5em, itemsep=10pt]
  \item Find $P(\text{Trains Daily} | \text{Swims})$. Show setup.\\
        \blank{3cm} $\div$ \blank{3cm} $=$ \blank{3cm}

  \item Find $P(\text{Cycles} | \text{Trains Weekly})$. Show setup.\\
        \blank{3cm} $\div$ \blank{3cm} $=$ \blank{3cm}

  \item Use the multiplication rule: $P(\text{Trains Weekly}) \cdot P(\text{Cycles}|\text{Trains Weekly}) = $ \blank{4cm}\\
        Does this match the cell count? \blank{2cm}

  \item Is $P(\text{Runs}|\text{Trains Daily})$ equal to $P(\text{Trains Daily}|\text{Runs})$?
        Calculate both and compare.\\
        \writeline\writeline
\end{enumerate}
\end{practicebox}

\end{document}
